% Source: physical PDF p. 229 (printed p. 206). Complete three-panel Figure 19.3 and caption.
\begin{figure}[htbp]
  \centering
  \begin{tikzpicture}[
    x=0.82cm,
    y=0.82cm,
    line cap=butt,
    line join=round,
    nucleon/.style={line width=0.9pt},
    pion/.style={line width=0.75pt,dash pattern=on 4pt off 4pt}
  ]
    % (a): direct nucleon pole graph.
    \draw[nucleon] (0,-2.05) -- (0,2.05);
    \draw[pion] (0,0.83) -- (2.10,2.08);
    \draw[pion] (0,-0.83) -- (2.10,-2.08);
    \node at (0,-2.85) {(a)};

    % (b): crossed nucleon pole graph.
    \begin{scope}[xshift=5.25cm]
      \draw[nucleon] (0,-2.05) -- (0,2.05);
      \draw[pion] (0,0.83) -- (2.10,-1.05);
      \draw[pion] (0,-0.83) -- (2.10,1.05);
      \node at (0,-2.85) {(b)};
    \end{scope}

    % (c): pion-pion-nucleon contact graph.
    \begin{scope}[xshift=10.50cm]
      \draw[nucleon] (0,-2.05) -- (0,2.05);
      \draw[pion] (0,0) -- (2.10,1.22);
      \draw[pion] (0,0) -- (2.10,-1.22);
      \node at (0,-2.85) {(c)};
    \end{scope}
  \end{tikzpicture}
  \renewcommand{\thefigure}{19.3}
  \caption{Feynman diagrams used with an effective chiral Lagrangian to calculate the scattering of a soft pion from a nucleon. Dashed lines are pions; solid lines are nucleons.}
  \label{fig:19.3}
\end{figure}
